\documentclass[letterpaper,10pt]{article}
\setlength{\footskip}{4.08003pt}


\usepackage{latexsym}
\usepackage[empty]{fullpage}
\usepackage{titlesec}
\usepackage[usenames,dvipsnames]{color}
\usepackage{enumitem}
\usepackage[hidelinks]{hyperref}
\usepackage{fancyhdr}
% \usepackage[english]{babel}
% \usepackage{tabularx}
% \usepackage{hyphenat}
\usepackage{fontawesome}

\pagestyle{fancy}
\fancyhf{} % clear all header and footer fields
\fancyfoot[L]{\small{Last Updated: \textit{\today} \ by \href{https://renzhenyu.site}{Zhenyu Ren}}} % 设置页脚内容
\renewcommand{\headrulewidth}{0pt}
\renewcommand{\footrulewidth}{0pt}

% Adjust margins
\addtolength{\oddsidemargin}{-0.5in}
\addtolength{\evensidemargin}{-0.5in}
\addtolength{\textwidth}{1in}
\addtolength{\topmargin}{-.5in}
\addtolength{\textheight}{1.0in}

\urlstyle{same}

\raggedbottom
\raggedright
\setlength{\tabcolsep}{0in}

% Sections formatting
\titleformat{\section}{
  \vspace{-4pt}\scshape\raggedright\large
}{}{0em}{}[\color{black}\titlerule \vspace{-5pt}]

% Ensure that generate pdf is machine readable/ATS parsable
\pdfgentounicode=1

%-------------------------
% Custom commands

\newcommand{\resumeItem}[1]{
  \item\small{
    {#1 \vspace{-2pt}}
  }
}


\newcommand{\resumeSubheading}[4]{
  \vspace{-2pt}\item
    \begin{tabular*}{0.97\textwidth}[t]{l@{\extracolsep{\fill}}r}
      \textbf{#1} & #2 \\
      \textit{\small#3} & \textit{\small #4} \\
    \end{tabular*}\vspace{-7pt}
}


\newcommand{\resumeSubSubheading}[2]{
    \vspace{-2pt}\item
    \begin{tabular*}{0.97\textwidth}{l@{\extracolsep{\fill}}r}
      \textit{\small#1} & \textit{\small #2} \\
    \end{tabular*}\vspace{-7pt}
}


\newcommand{\resumeEducationHeading}[6]{
  \vspace{-2pt}\item
    \begin{tabular*}{0.97\textwidth}[t]{l@{\extracolsep{\fill}}r}
      \textbf{#1} & #2 \\
      \textit{\small#3} & \textit{\small #4} \\
      \textit{\small#5} & \textit{\small #6} \\
    \end{tabular*}\vspace{-5pt}
}


\newcommand{\resumeProjectHeading}[2]{
    \vspace{-2pt}\item
    \begin{tabular*}{0.97\textwidth}{l@{\extracolsep{\fill}}r}
      \small#1 & #2 \\
    \end{tabular*}\vspace{-7pt}
}


\newcommand{\resumeOrganizationHeading}[4]{
  \vspace{-2pt}\item
    \begin{tabular*}{0.97\textwidth}[t]{l@{\extracolsep{\fill}}r}
      \textbf{#1} & \textit{\small #2} \\
      \textit{\small#3}
    \end{tabular*}\vspace{-7pt}
}

\newcommand{\resumeSubItem}[1]{\resumeItem{#1}\vspace{-4pt}}

\renewcommand\labelitemii{$\vcenter{\hbox{\tiny$\bullet$}}$}

\newcommand{\resumeSubHeadingListStart}{\begin{itemize}[leftmargin=0.15in, label={}]}
\newcommand{\resumeSubHeadingListEnd}{\end{itemize}\vspace{-5pt}}
\newcommand{\resumeItemListStart}{\begin{itemize}}
\newcommand{\resumeItemListEnd}{\end{itemize}\vspace{-5pt}}


% Quotation environment (for reply)
\newenvironment{myQuote}{
    \bigskip
    \begin{mdframed}
        [topline=false,
            rightline=false,
            bottomline=false,
            innertopmargin=0.4em,
            innerbottommargin=0.4em,
            innerrightmargin=0.7em,
            rightmargin=0.7em,
            innerleftmargin=0.7em,
            leftmargin=0.7em,
            linewidth=.2em,
            linecolor=black]
        % Handling footnote style in mdframed
        \stepcounter{footnote}
        \renewcommand{\thempfootnote}{\arabic{footnote}}
        }{
    \end{mdframed}
    \bigskip
}


%-------------------------------------------
%%%%%%  RESUME STARTS HERE  %%%%%%%%%%%%%%%%%%%%%%%%%%%%


\begin{document}

%---------- HEADING ----------

\begin{center}
    \textbf{\Huge \scshape {Zhenyu Ren}} \\ \vspace{3pt}
    \small
    % \faMobile \hspace{.5pt} \href{tel:905314204536}{+90 531 420 4536}
    % $|$
    % \faAt \hspace{.5pt} Mail: \href{mailto:renzy2022@mail.sustech.edu.cn}{\textcolor{blue}{renzy2022@mail.sustech.edu.cn}}
    \faAt \hspace{.5pt} Mail: \href{mailto:zrenax@connect.ust.hk}{\textcolor{blue}{zrenax@connect.ust.hk}}
    $|$
    % \faLinkedinSquare \hspace{.5pt} \href{https://www.linkedin.com/in/arasgungore}{LinkedIn}
    % $|$
    \faGithub \hspace{.5pt} Github: \href{https://github.com/rzy0901}{\textcolor{blue}{https://github.com/rzy0901}}
    $|$
    \faGlobe \hspace{.5pt} Website: \href{https://renzhenyu.site}{\textcolor{blue}{https://renzhenyu.site}}
    % $|$
    % \faMapMarker \hspace{.5pt} {Shenzhen, China}
\end{center}



%----------- EDUCATION -----------

\section{Education}
  \vspace{3pt}
  \resumeSubHeadingListStart
      \resumeSubheading
      {Hong Kong University of Science and Technology (HKUST)}{Hong Kong, China}
      {PhD student in Computer Science and Engineering}{Sep 2025 \textbf{--} Present}
    \resumeSubHeadingListStart
        \small{\item{Supervised by Prof.~\href{https://home.cse.ust.hk/~lim/}{\textcolor{blue}{Mo~Li}} (IEEE Fellow), focusing on short range wireless communication protocols and security (e.g., Nearlink, BLE, WiFi).}}
    \resumeSubHeadingListEnd 
    \resumeEducationHeading
      {Southern University of Science and Technology (SUSTech)
      }{Shenzhen, China}
      {B.Sc. in Communication Engineering;   \textbf{GPA: 3.69/4.00}; \textbf{Weighted Score: 88.07}}{Sep 2018 \textbf{--} Jun 2022}
      {M.Sc. in Electronic Science and Technology;
      \textbf{GPA: 3.44/4.00}; \textbf{Weighted Score: 88.52} }{Sep 2022 \textbf{--} Jun 2025}
        \resumeSubHeadingListStart
        % \small{\item{Supervised by Prof. \href{https://ieeexplore.ieee.org/author/37089502484}{\textcolor{blue}{Rui Wang}} (Area Editor of IEEE OJ-COMS), focusing on \textbf{sensing channel modeling and experimental mmwave system for wireless sim2real gesture recognition.} \\ \vspace{3pt}
        \small{\item{Supervised by Prof. \href{https://ieeexplore.ieee.org/author/37089502484}{\textcolor{blue}{Rui Wang}} (Area Editor of IEEE OJ-COMS), focusing on wireless sensing. \\ \vspace{3pt}
        }}
        \resumeSubHeadingListEnd    
  \resumeSubHeadingListEnd

%----------- EXPERIENCE -----------

\section{Experience}
  \vspace{3pt}
  \resumeSubHeadingListStart
    \resumeSubheading
      {Huawei Hong Kong Research Center (Leibniz Institute)}{Hong Kong, China}
      {Part-time Research Intern}{June 2026 \textbf{--} Present}
      \resumeSubHeadingListStart
      \small{\item{Working on the \textbf{SparkLink Low Energy (SLE) protocol} for multi-device collaboration in HarmonyOS.}
      }
      \resumeSubHeadingListEnd
    \resumeSubheading
      {Huawei Technology, Computer Network and Protocol Research Laboratory}{Shenzhen, China}
      {Research Intern}{July 2024 \textbf{--} Jan 2025}
      \resumeSubHeadingListStart
      \small{\item{Working on \textbf{Nearlink (Sparklink) wireless communication protocal}. 
      % Collaborating with Prof.~\href{https://home.cse.ust.hk/~lim/}{\textcolor{blue}{Mo~Li}} (IEEE Fellow).
      }
      }
      \resumeSubHeadingListEnd  
    \resumeSubheading
      {Huawei Technology, Wireless Technology (WT) Laboratory}{Shenzhen, China}
      {Research Intern}{June 2021 \textbf{--} June 2022}
      \resumeSubHeadingListStart
      \small{\item{Focusing on \textbf{wireless communication \& sensing technology}. 
      % Supervised by Dr. \href{https://ieeexplore.ieee.org/author/37088772113}{\textcolor{blue}{Tony Han Xiao}} (Chair of IEEE 802.11bf).
      }
      }
      \resumeSubHeadingListEnd  
  \resumeSubHeadingListEnd

%----------- PUBLICATIONS -----------
\section{Publications}
\vspace{3pt}
\resumeSubHeadingListStart
\small{\item{
    \textbf{Zhenyu Ren}, Yanbo Zhang, Boya Liu, and Mo Li. ``Snatcher: Apple Find My Network Exposes Your Lost Devices To Strangers," in \textit{Proceedings of the 2026 ACM SIGSAC Conference on Computer and Communications Security (CCS)}, 2026.
    \href{https://arxiv.org/abs/2606.21067}{\color{blue}Paper} $|$ \href{https://github.com/rzy0901/Snatcher}{\color{blue}Artifact}.
    Awarded \textbf{Artifacts Available}, \textbf{Artifacts Evaluated -- Functional}, and \textbf{Results Reproduced} badges.

    Boya Liu*, \textbf{Zhenyu Ren*}, Yanbo Zhang, Mo Li, Jiaxin Liang, Bin He, Jie Li, Hao Wu, and Jingbin Zhou. 2026. Unveiling SparkLink Low Energy: A Comparative Measurement Study. \textit{ACM Trans. Internet Things} 7, 1, Article 6 (February 2026), 24 pages. https://doi.org/10.1145/3776559. (Co-first author).

    \textbf{Zhenyu Ren}, Guoliang Li, Chenqing Ji, Chao Yu, Shuai Wang, and Rui Wang. ``CASTER: A Computer-Vision-Assisted Wireless Channel Simulator for Gesture Recognition," in \textit{IEEE Open Journal of the Communications Society}, vol. 5, pp. 3185-3195, 2024. DOI: 10.1109/OJCOMS.2024.3398016.
    \href{https://doi.org/10.1109/ojcoms.2024.3398016}{\color{blue}Paper} $|$ \href{https://github.com/rzy0901/testSpectrogram}{\color{blue}GitHub}  $|$ \href{https://lasso-sustech.github.io/CASTER/}{\color{blue}Project Page}.
    }
    \item{\textbf{REN Zhenyu}, JI Chenqing, YU Chao, CHEN Wanli, and WANG Rui. Computer vision-assisted wireless channel simulation for millimeter wave human motion recognition[J]. \textit{Journal of Radars},  in press. DOI:  10.12000/JR24101.
    }
}
\resumeSubHeadingListEnd

% %----------- PATENTS -----------
% \section{Patents}
% \vspace{3pt}
% % 任振裕，陈万里，王锐，余潮，“无线信道仿真方法、装置、计算机设备及存储介质”，专利申请号：2023110356420，南方科技大学，申请日：2023.08.16（中国发明专利）。
% % 翻译上述专利
% \resumeSubHeadingListStart
% \small{\item{
%       \textbf{Zhenyu Ren}, Wanli Chen, Rui Wang, Chao Yu, “Wireless Channel Simulation Method, Device, Computer Equipment, and Storage Medium”, Patent Application No.: 2023110356420, Southern University of Science and Technology, Application Date: 2023.08.16 (Chinese Invention Patent).
%     }
% }

% \resumeSubHeadingListEnd
%----------- AWARDS & ACHIEVEMENTS -----------

\section{Skills}
\resumeSubHeadingListStart
    \item Claude, Codex, Gemini, Kimi, Cursor, GitHub Copilot, ChatGPT, OpenClaw...
\resumeSubHeadingListEnd

% \section{Awards \& Achievements}
%   \vspace{3pt}
%   \resumeSubHeadingListStart
%     \small{\item{
%       Second Prize in the 17th ``Challenge Cup" Guangdong University Student Extracurricular Academic Science and Technology Works Competition, 2023.
%       \\ \vspace{3pt}
%       Leader for Guangdong University Students' Science and Technology Innovation Cultivation Special Fund (``Climbing Plan" Special Fund), 2022$\sim$2023 \emph{(Funding: 20,000 RMB)}.
%       \\ \vspace{3pt}
%       2022 Excellent Graduate of Undergraduate for exceptional performance in the SUSTech.
%     %   \\ \vspace{3pt}
%     %   2022 Distinguished Undergraduate Thesis of the SUSTech.
%       \\ \vspace{3pt}
%       Southern University of Science and Technology Outstanding Student Third-Class Scholarship (2018$\sim$2019, 2020$\sim$2021).
%       \\ \vspace{3pt}
%       First Prize in the 2020 National College Student Mathematics Modeling Competition \emph{(Top 0.5\% among 45000++ teams in China)}.
%     }}
%   \resumeSubHeadingListEnd

%----------- SKILLS -----------

% \section{Skills}
%   \vspace{3pt}
%   \resumeSubHeadingListStart
%     \small{\item{
%         \textbf{Programming Languages:}{ C/C++, Python, MATLAB, Java} \\ \vspace{3pt}
%         \textbf{Technologies:}{ PyTorch, Linux/Ubuntu, Git/GitHub, OpenCV, UHD/USRP, 60GHz Sivers}
%         \\ \vspace{3pt}
%         \textbf{Writing:}{ \LaTeX, Markdown, Website} \\ \vspace{3pt}
%         \textbf{English:}{ IELTS 6.5 with sub-score 6.0} (Test date: \textit{\small Jan 2024});
%     }}
%   \resumeSubHeadingListEnd

% %----------- PROJECTS -----------

% \section{Selected Projects}
% \vspace{3pt}
% \emph{Note: You can look up my \href{https://github.com/rzy0901}{\textcolor{blue}{GitHub}} for my full project lists.}
%     \resumeSubHeadingListStart
%       \resumeProjectHeading
%         {\textbf{CASTER} $|$           \href{https://doi.org/10.1109/ojcoms.2024.3398016}{\color{blue}Paper} $|$ \href{https://github.com/rzy0901/testSpectrogram}{\color{blue}GitHub} $|$ \href{https://lasso-sustech.github.io/CASTER/}{\textcolor{blue}{Project Page}}}{}
%         %   \resumeItemListStart
%             % \resumeItem{CASTER: A {C}omputer-vision-{A}ssisted wireless channel {S}imula{T}or for g{E}sture {R}ecognition.}
%             \resumeItem{An open-source platform for wireless channel simulation, human/hand pose extraction, gesture spectrogram generation, and real-time gesture recognition based on millimeter-wave passive sensing and communication systems.
%             % \resumeItemListStart
%             %   \resumeItem{Submodules \href{https://github.com/rzy0901/mediapipe_spectrogram}{\textcolor{blue}{mediapipe\_spectrogram}} and \href{https://github.com/rzy0901/testZED}{\textcolor{blue}{testZED}}: Developed algorithms for keypoint extraction from video streams and used a primitive-based channel model to generate simulated data, addressing the data collection issue in wireless sensing.}
%             %   \resumeItem{Submodule \href{https://github.com/rzy0901/CASTER_classification}{\textcolor{blue}{CASTER\_classification}}: Implemented a Simulation-to-Reality transfer learning strategy using ResNet18 and adversarial discriminative domain adaptation (ADDA) for wireless gesture recognition. This approach improved real-world dataset accuracy from $83.0\%$ to $96.5\%$.}
%             %   \resumeItem{Submodule \href{https://github.com/rzy0901/RxRealTime_GUI_rzy}{\textcolor{blue}{RxRealTime\_GUI}}: Implemented real-time gesture recognition based on millimeter-wave passive sensing and communication systems, using USRP and 60GHz Sivers phased array.
%             %   }
%             % \resumeItemListEnd
%             }         
%         %   \resumeItemListEnd
%           \small{\item{\textbf{Some other projects:}\\ \vspace{3pt} (1) \href{https://github.com/rzy0901/rayTrace}{\textcolor{blue}{rayTrace}}: MATLAB software for indoor radio ray-tracing with human blockage. \textbf{\emph{A segment of GIF results was featured by Huawei's keynote speaker during the 1st workshop on wi-fi sensing (\href{https://www.bilibili.com/video/BV1WK411C7hB?t=951.8}{\textcolor{blue}{video recording at 15:52}}).}}\vspace{3pt} \\
%           (2) \href{https://github.com/rzy0901/rzy0901.github.io}{\textcolor{blue}{Personal Website}}: Math notes and technical blogs.}}
%         %   \small{\item{\textbf{Some other projects:}\\ \vspace{3pt} (1) \href{https://rzy0901.github.io/testWinprop/}{\textcolor{blue}{testWinprop}}: Radio Ray-tracing by Winprop C++ API (7 Stars); (2) \href{https://github.com/rzy0901/My-solutions-to-DSP-LAB}{\textcolor{blue}{My-solutions-to-DSP-LAB}}: MATLAB for Digital Signal Processing Labs (16 stars, 1 fork); (3) \href{https://github.com/rzy0901/CS205_Matrix}{\textcolor{blue}{CS205\_Matrix}}: C++ Matrix computation library (6 stars, 2 forks); (4) \href{https://github.com/rzy0901/rzy0901.github.io}{\textcolor{blue}{Personal Website}}: Math notes and technical blogs (example post: \href{https://renzhenyu.site/post/asp/}{\textcolor{blue}{Notebook for Applied Stochastic Processes}} and \href{https://renzhenyu.site/post/channel/}{\textcolor{blue}{Wireless Channel Trouble-Shooting}}).}}  
%     \resumeSubHeadingListEnd
    
    
% %----------- Extracurricular Activities -----------
% \section{Extracurricular Activities}
% \vspace{3pt}
% \resumeSubHeadingListStart
% \small{
%   \item{
%   Reviewer for IEEE ICC 2024, IEEE GLOBECOM 2024, IEEE ICMLCN 2024.\\ \vspace{3pt}
%   TPC reviewer for IEEE ICC 2024 Workshop. \\ \vspace{3pt}
%   Teaching Assistant for Wireless Communication.
% }}
% \resumeSubHeadingListEnd

%----------- REFERENCES -----------
% \section{References}
% \vspace{3pt}
% \resumeSubHeadingListStart
% \small{
%   \item{
%   Prof. \href{https://ieeexplore.ieee.org/author/37089502484}{\textcolor{blue}{Rui Wang}}, Associate Professor, Department of Electronic and Electrical Engineering (EEE), Southern University of Science and Technology, Email: \href{mailto:wang.r@sustech.edu.cn}{\textcolor{blue}{wang.r@sustech.edu.cn}}.}}
% \resumeSubHeadingListEnd

% %----------- FOOTER -----------
%   % \vspace{\baselineskip}
%   \titlerule
%   \begin{center}
%       \small{
%       Last Updated: \textit{July 07, 2024} \ by \href{https://renzhenyu.site}{Zhenyu Ren}}
%   \end{center}

\end{document}
